%! Introduction to Eigenvalues — Unit 6, Lesson 6.1 — Warm-Up
\documentclass[10pt]{article}
\usepackage{linalg-article}
\usepackage{linalg-boxes}
\newcommand{\TallMath}[1]{$\displaystyle #1\rule[-1.4em]{0pt}{3.2em}$}

\begin{document}

\pageheader{Unit 6, Lesson 6.1}{Warm-Up}
\namedateperiod

\vspace{0.12in}

\begin{notesbox}{Getting Ready: The Three Moves We'll Use Today}
\small Yesterday you \emph{tested} vectors you were handed. Today we \emph{find} the special
directions ourselves --- and these three moves are the whole method. Answer quickly.

\vspace{4pt}
\textbf{1. Subtract $\lambda$ down the diagonal.} For $A=\begin{bmatrix} 4 & 1 \\ 2 & 3 \end{bmatrix}$,
write the matrix $A-\lambda I$ (subtract the unknown $\lambda$ from each diagonal entry):
\[
  A-\lambda I=\begin{bmatrix}\rule{0pt}{1.2em}\blank{1.4cm} & \blank{1.0cm}\\[3pt]\blank{1.0cm} & \blank{1.4cm}\end{bmatrix}.
\]

\vspace{4pt}
\textbf{2. The $2\times2$ determinant (Unit 5).} Recall $\det\begin{bsmallmatrix} a & b \\ c & d
\end{bsmallmatrix}=ad-bc$. Compute
\[
  \det\begin{bmatrix} 4 & 1 \\ 2 & 3 \end{bmatrix} = \blank{2.4cm}.
\]

\vspace{4pt}
\textbf{3. Solve a quadratic (factor).} Find both values of $\lambda$ that make this true:
\[
  \lambda^2-5\lambda+6=0 \qquad\Longrightarrow\qquad \lambda=\blank{1.2cm}\ \text{ or }\ \lambda=\blank{1.2cm}.
\]
\end{notesbox}

\end{document}
